\documentclass[tikz]{standalone} % Standalone: Only graphics export.
\usepackage[european,siunitx]{circuitikz}
\usetikzlibrary{arrows.meta}
\usepackage{listings}
\lstset{basicstyle=\footnotesize\ttfamily}
\begin{document}

  \newcommand{\eq}{=}
  \tikzset{ameter/.style={rmeterwa, t=A}}
  \begin{circuitikz}

    \tikzset{>=Triangle} % Define tip of arrows drawn with [->]
    \ctikzset{voltage=straight}


    \def\Lr{2.5}
    \def\deltaUinY{1}
    \def\deltaUinX{1.2}
    \def\sizeY{3}
    \def\deltaUio{1}
    \def\gapUin{0.3}
    
    % Horizontal upper line.
    \draw (0,\sizeY) coordinate(in1)
      to[short,o-] ++(\deltaUinX,0)
      to[R,l=$R_1$,-*] ++(\Lr,0) coordinate(Kl) node[above]{$K_l$}
      to[L,l=$L$] ++(\Lr,0) coordinate(Klout) %node[above]{$K_{l_{out}}$} 
      to[ameter,i=$I_L$,v=$U_{probe}$,-*] ++(\Lr,0) coordinate(TopC)
      to[short,-o] ++(\deltaUio,0) coordinate(out) node[above]{$out$};
    
    % Second input line.
    \draw (\deltaUinX,{\sizeY-\deltaUinY}) coordinate(in2)
      to[R,l=$R_2$,o-] ++(\Lr,0)
      -- (Kl);

    % Down to C and bottom output connector.
    \draw (TopC)
      to[C,l_=$C$,-*] ++(0,-\sizeY) coordinate(BottomC)
      to[short,-o] ++(\deltaUio,0) coordinate(BottomUout);

    % Horizontal bottom (ground) line.
    \draw (0,0) coordinate(BottomUin1)
      to[short,o-o] ++(\deltaUinX,0) coordinate(BottomUin2)
      -- (BottomC);

    % Ground node.
    \draw (BottomC)
      -- ++(0,-0.4) node[ground]{};

    % Voltage arrow for first input voltage.
    \draw[->] ($(in1)+(0,-\gapUin)$)
      -- node[midway,left]{$U_{in_1}$} ($(BottomUin1)+(0,\gapUin)$);

    % Voltage arrow for second input voltage.
    \draw[->] ($(in2)+(0,-\gapUin)$)
      -- node[midway,left]{$U_{in_2}$} ($(BottomUin2)+(0,\gapUin)$);

    % Voltage arrow for output voltage.
    \draw[->] ($(out)+(0,-\gapUin)$)
      -- node[midway,right]{$U_{out}$} ($(BottomUout)+(0,\gapUin)$);

    % Voltage arrow for voltage U_LC.
    \draw[->] ($(in2)+(\Lr,0)+(0,-\gapUin)$)
      -- node[midway,left]{$U_{LC}$} ++(0,{-\sizeY+\deltaUinY+2*\gapUin});

  \end{circuitikz}
\end{document}
